% ==============================================================================
% SLIDES - SEMANA 13: QUALIDADE E MATURIDADE (CMMI, MPS-SW, ISO 12207, ISO 25010)
% ==============================================================================

% 1. CAPA
\senaicover
  {Qualidade \& Maturidade de Processos}
  {CMMI-DEV, MR-MPS-SW, ISO/IEC 12207, ISO 25010 e Entrega Parcial 3}
  {Unidade Curricular: Governança de TI}
  {Prof. André Souza}

% 2. AGENDA
\begin{senaiframe}{Visão Geral \& Agenda}{Estrutura da Aula}
  \begin{columns}[t]
    \begin{column}{0.48\textwidth}
      \begin{tcolorbox}[senaibluecard, title={1. Modelos de Maturidade}]
        \begin{itemize}
          \item CMMI-DEV (Níveis 1 a 5).
          \item MR-MPS-SW (Níveis G a A).
        \end{itemize}
      \end{tcolorbox}
      \vspace{4pt}
      \begin{tcolorbox}[senaiambercard, title={3. Qualidade de Produto}]
        \begin{itemize}
          \item ISO/IEC 25010 (SQuaRE).
          \item As 8 características de qualidade de software.
        \end{itemize}
      \end{tcolorbox}
    \end{column}
    \begin{column}{0.48\textwidth}
      \begin{tcolorbox}[senairedcard, title={2. Ciclo de Vida (ISO 12207)}]
        \begin{itemize}
          \item Processos Primários, Apoio e Organizacionais.
          \item Gestão de Mudanças e Configuração.
        \end{itemize}
      \end{tcolorbox}
      \vspace{4pt}
      \begin{tcolorbox}[senaigreencard, title={4. Entrega Parcial 3}]
        \begin{itemize}
          \item Submissão dos Capítulos 7 e 8.
          \item Terceiro marco do projeto PDGTI.
        \end{itemize}
      \end{tcolorbox}
    \end{column}
  \end{columns}
\end{senaiframe}

% 3. PROCESSO VS PRODUTO
\begin{senaiframe}{Teoria • Módulo 1}{Qualidade do Processo vs. Qualidade do Produto}
  \begin{columns}[t]
    \begin{column}{0.48\textwidth}
      \begin{tcolorbox}[senaibluecard, title={Premissa Fundamental}]
        \begin{itemize}
          \item A qualidade do software é resultado direto da qualidade do processo.
          \item Processos caóticos geram sistemas instáveis e cheios de bugs.
        \end{itemize}
      \end{tcolorbox}
    \end{column}
    \begin{column}{0.48\textwidth}
      \begin{tcolorbox}[senaigreencard, title={Foco em Processos}]
        \begin{itemize}
          \item Padronização, testes automatizados e controle de mudanças.
        \end{itemize}
      \end{tcolorbox}
    \end{column}
  \end{columns}
\end{senaiframe}

% 4. MODELOS DE MATURIDADE
\begin{senaiframe}{Teoria • Módulo 1}{Visão Geral dos Modelos de Maturidade}
  \begin{columns}[t]
    \begin{column}{0.48\textwidth}
      \begin{tcolorbox}[senaibluecard, title={O que são?}]
        \begin{itemize}
          \item Roteiros organizados em níveis para evolução de processos.
          \item Medem a capacidade da empresa em entregar código sem falhas.
        \end{itemize}
      \end{tcolorbox}
    \end{column}
    \begin{column}{0.48\textwidth}
      \begin{tcolorbox}[senaicard, title={Principais Referências}]
        \begin{itemize}
          \item CMMI-DEV (Internacional) e MR-MPS-SW (Nacional / Softex).
        \end{itemize}
      \end{tcolorbox}
    \end{column}
  \end{columns}
\end{senaiframe}

% 5. CMMI-DEV V2.0
\begin{senaiframe}{Teoria • Módulo 1}{CMMI-DEV v2.0 (5 Níveis de Maturidade)}
  \begin{columns}[t]
    \begin{column}{0.48\textwidth}
      \begin{tcolorbox}[senairedcard, title={Níveis Inicial a Definido}]
        \begin{itemize}
          \item \textbf{Nível 1 (Inicial):} Caótico, dependente de heróis.
          \item \textbf{Nível 2 (Gerenciado):} Gestão por projetos.
          \item \textbf{Nível 3 (Definido):} Processos padronizados.
        \end{itemize}
      \end{tcolorbox}
    \end{column}
    \begin{column}{0.48\textwidth}
      \begin{tcolorbox}[senaigreencard, title={Níveis Avançados}]
        \begin{itemize}
          \item \textbf{Nível 4 (Gerenciado Quantitativamente):} Métricas.
          \item \textbf{Nível 5 (Em Otimização):} Melhoria contínua.
        \end{itemize}
      \end{tcolorbox}
    \end{column}
  \end{columns}
\end{senaiframe}

% 6. MR-MPS-SW
\begin{senaiframe}{Teoria • Módulo 1}{O Modelo Brasileiro MR-MPS-SW (Softex)}
  \begin{columns}[t]
    \begin{column}{0.48\textwidth}
      \begin{tcolorbox}[senaiambercard, title={Sete Níveis de Maturidade}]
        \begin{itemize}
          \item Do Nível G (Parcialmente Gerenciado) até o Nível A (Em Otimização).
          \item Modelo mais acessível e de custo menor para PMEs.
        \end{itemize}
      \end{tcolorbox}
    \end{column}
    \begin{column}{0.48\textwidth}
      \begin{tcolorbox}[senaibluecard, title={Reconhecimento}]
        \begin{itemize}
          \item Totalmente compatível com CMMI e ISO/IEC 12207.
        \end{itemize}
      \end{tcolorbox}
    \end{column}
  \end{columns}
\end{senaiframe}

% 7. EQUIVALÊNCIA CMMI VS MPS-SW
\begin{senaiframe}{Teoria • Módulo 1}{Equivalência: CMMI vs. MR-MPS-SW}
  \begin{columns}[t]
    \begin{column}{0.48\textwidth}
      \begin{tcolorbox}[senaibluecard, title={CMMI}]
        \begin{itemize}
          \item 5 Níveis por Estágios.
          \item Custo elevado de certificação formal internacional.
        \end{itemize}
      \end{tcolorbox}
    \end{column}
    \begin{column}{0.48\textwidth}
      \begin{tcolorbox}[senaigreencard, title={MPS-SW}]
        \begin{itemize}
          \item 7 Níveis (G a A) permitindo implantação em degraus menores.
          \item Foco no mercado nacional e editais públicos.
        \end{itemize}
      \end{tcolorbox}
    \end{column}
  \end{columns}
\end{senaiframe}

% 8. ISO/IEC 12207
\begin{senaiframe}{Teoria • Módulo 2}{A Norma ISO/IEC 12207}
  \begin{columns}[t]
    \begin{column}{0.48\textwidth}
      \begin{tcolorbox}[senaibluecard, title={Ciclo de Vida de Software}]
        \begin{itemize}
          \item Estrutura comum para desenvolvimento, operação e manutenção.
          \item Define a arquitetura de processos da engenharia de software.
        \end{itemize}
      \end{tcolorbox}
    \end{column}
    \begin{column}{0.48\textwidth}
      \begin{tcolorbox}[senaicard, title={Padronização}]
        \begin{itemize}
          \item Base conceitual utilizada pelo CMMI e MPS-SW.
        \end{itemize}
      \end{tcolorbox}
    \end{column}
  \end{columns}
\end{senaiframe}

% 9. GRUPOS DA ISO 12207
\begin{senaiframe}{Teoria • Módulo 2}{Processos da ISO/IEC 12207}
  \begin{columns}[t]
    \begin{column}{0.48\textwidth}
      \begin{tcolorbox}[senaibluecard, title={Primários \& Apoio}]
        \begin{itemize}
          \item \textbf{Primários:} Aquisição, Desenvolvimento, Operação.
          \item \textbf{Apoio:} Garantia da Qualidade, Auditoria, Mudanças.
        \end{itemize}
      \end{tcolorbox}
    \end{column}
    \begin{column}{0.48\textwidth}
      \begin{tcolorbox}[senaiambercard, title={Organizacionais}]
        \begin{itemize}
          \item \textbf{Organizacionais:} Gestão de Processos e Treinamento.
        \end{itemize}
      \end{tcolorbox}
    \end{column}
  \end{columns}
\end{senaiframe}

% 10. QUALIDADE DO PRODUTO (ISO 25010)
\begin{senaiframe}{Teoria • Módulo 3}{ISO/IEC 25010 — Qualidade do Produto}
  \begin{columns}[t]
    \begin{column}{0.48\textwidth}
      \begin{tcolorbox}[senaibluecard, title={Evolução da ISO 9126}]
        \begin{itemize}
          \item Norma SQuaRE para avaliação da qualidade de produto de software.
          \item Define 8 características essenciais auditáveis.
        \end{itemize}
      \end{tcolorbox}
    \end{column}
    \begin{column}{0.48\textwidth}
      \begin{tcolorbox}[senaigreencard, title={Objetivo}]
        \begin{itemize}
          \item Avaliação objetiva do software antes da entrega final ao cliente.
        \end{itemize}
      \end{tcolorbox}
    \end{column}
  \end{columns}
\end{senaiframe}

% 11. CARACTERÍSTICAS DA ISO 25010 - PARTE 1
\begin{senaiframe}{Teoria • Módulo 3}{ISO 25010: Operação \& Segurança}
  \begin{columns}[t]
    \begin{column}{0.48\textwidth}
      \begin{tcolorbox}[senairedcard, title={Adequação \& Eficiência}]
        \begin{itemize}
          \item \textbf{Adequação Funcional:} Precisão nos requisitos.
          \item \textbf{Eficiência de Desempenho:} Tempo de resposta.
        \end{itemize}
      \end{tcolorbox}
    \end{column}
    \begin{column}{0.48\textwidth}
      \begin{tcolorbox}[senaiambercard, title={Segurança \& Compatibilidade}]
        \begin{itemize}
          \item \textbf{Segurança:} Proteção contra acessos indevidos.
          \item \textbf{Compatibilidade:} Interoperabilidade entre sistemas.
        \end{itemize}
      \end{tcolorbox}
    \end{column}
  \end{columns}
\end{senaiframe}

% 12. CARACTERÍSTICAS DA ISO 25010 - PARTE 2
\begin{senaiframe}{Teoria • Módulo 3}{ISO 25010: Usabilidade \& Manutenção}
  \begin{columns}[t]
    \begin{column}{0.48\textwidth}
      \begin{tcolorbox}[senaibluecard, title={Usabilidade \& Confiabilidade}]
        \begin{itemize}
          \item \textbf{Usabilidade:} Aprendizado simples e acessibilidade.
          \item \textbf{Confiabilidade:} Tolerância a falhas e resiliência.
        \end{itemize}
      \end{tcolorbox}
    \end{column}
    \begin{column}{0.48\textwidth}
      \begin{tcolorbox}[senaigreencard, title={Manutenção \& Portabilidade}]
        \begin{itemize}
          \item \textbf{Manutenibilidade:} Código limpo e testável.
          \item \textbf{Portabilidade:} Funcionamento em múltiplos ambientes.
        \end{itemize}
      \end{tcolorbox}
    \end{column}
  \end{columns}
\end{senaiframe}

% 13. ESTUDO DE CASO PRÁTICO
\begin{senaiframe}{Estudo de Caso}{Software House: Da Crise ao MPS-SW Nível F}
  \begin{columns}[t]
    \begin{column}{0.48\textwidth}
      \begin{tcolorbox}[senairedcard, title={Cenário Inicial}]
        \begin{itemize}
          \item Entregas com atraso, bugs constantes e insatisfação de clientes.
          \item Desenvolvimento informal sem controle de versão.
        \end{itemize}
      \end{tcolorbox}
    \end{column}
    \begin{column}{0.48\textwidth}
      \begin{tcolorbox}[senaigreencard, title={Evolução com MPS-SW}]
        \begin{itemize}
          \item Implantação do Nível G (Gerenciamento de Requisitos e Projetos).
          \item Redução de 70\% nos defeitos em produção após 1 ano.
        \end{itemize}
      \end{tcolorbox}
    \end{column}
  \end{columns}
\end{senaiframe}

% 14. REFERÊNCIAS BIBLIOGRÁFICAS
\begin{senaiframe}{Referências Bibliográficas}{Fontes \& Leitura Recomendada}
  \begin{tcolorbox}[senaicard, title={Obras de Referência}]
    \begin{itemize}
      \item \textbf{SOFTEX.} *MPS.BR - Melhoria de Processo do Software Brasileiro: Guia Geral.* Softex, 2021.
      \item \textbf{CMMI INSTITUTE.} *CMMI for Development, Version 2.0.* ISACA, 2018.
      \item \textbf{ISO/IEC.} *ISO/IEC 25010: Systems and software engineering — SQuaRE.* ISO, 2011.
    \end{itemize}
  \end{tcolorbox}
\end{senaiframe}

% 15. APLICAÇÃO NO PDGTI
\begin{senaiframe}{Aplicação no PDGTI}{Orientações para a ENTREGA PARCIAL 3}
  \vspace{0.2cm}
  \begin{tcolorbox}[senaigreencard, title={Entrega Parcial 3 (Capítulos 7 e 8) – Data: 30/Out}]
    \begin{itemize}
      \item \textbf{Capítulo 7:} Seleção e Justificativa de Frameworks (ISO 38500, COBIT, ITIL).
      \item \textbf{Capítulo 8:} Avaliação CMMI/MPS-SW, Métricas ISO 25010 e Mudanças.
      \item Consolidar com as correções dos Capítulos 1 a 6.
    \end{itemize}
  \end{tcolorbox}
\end{senaiframe}
